\documentclass[11pt,a4paper]{article}

\usepackage[utf8]{inputenc}
\usepackage[T1]{fontenc}
\usepackage[a4paper,margin=2.2cm]{geometry}
\usepackage{xcolor}
\usepackage{hyperref}
\usepackage{longtable}
\usepackage{array}
\usepackage{listings}

\hypersetup{
  colorlinks=true,
  linkcolor=blue,
  urlcolor=blue
}

\definecolor{codebg}{RGB}{245,247,250}
\definecolor{codeframe}{RGB}{220,225,232}

\lstdefinestyle{evalon}{
  basicstyle=\ttfamily\small,
  backgroundcolor=\color{codebg},
  frame=single,
  rulecolor=\color{codeframe},
  breaklines=true,
  columns=fullflexible,
  keepspaces=true
}

\title{Evalon Graph WebView Cloud Run Kullanim Dokumani}
\author{Codex}
\date{17 Nisan 2026}

\begin{document}
\maketitle

\section{Ozet}

Bu dokuman, \texttt{backend/graph} icindeki chart istemcisinin Cloud Run uzerine nasil alindigini ve web ile mobile tarafinda bu sayfanin nasil WebView olarak kullanilacagini aciklar.

Mevcut durum:

\begin{itemize}
  \item Chart istemcisi artik ayri bir Vite dev server'a bagimli degildir.
  \item Ayni Cloud Run hostu hem static sayfalari hem de \texttt{/api/v1/*} endpointlerini servis eder.
  \item Bu servis, veri icin \texttt{evalon-backtest-api} Cloud Run backend'ine gider.
  \item Web veya mobile uygulama yalnizca dogru URL'i olusturup WebView acarak chart'i gosterebilir.
\end{itemize}

\section{Canli Production Host}

Ana host:

\begin{lstlisting}[style=evalon]
https://evalon-graph-web-474112640179.europe-west1.run.app
\end{lstlisting}

Chart icin onerilen URL:

\begin{lstlisting}[style=evalon]
https://evalon-graph-web-474112640179.europe-west1.run.app/chart?symbol=ADEL&tf=1d
\end{lstlisting}

Alternatif tam dosya path'i:

\begin{lstlisting}[style=evalon]
https://evalon-graph-web-474112640179.europe-west1.run.app/chart.html?symbol=ADEL&tf=1d
\end{lstlisting}

Ek sayfalar:

\begin{itemize}
  \item \texttt{/backtest} veya \texttt{/backtest.html}
  \item \texttt{/ai} veya \texttt{/ai.html}
\end{itemize}

\section{Mimari}

Deploy edilen servis tek host uzerinden su akisi izler:

\begin{enumerate}
  \item Kullanici uygulamada hisse secer.
  \item Uygulama \texttt{symbol} ve \texttt{tf} query parametreleriyle chart URL'ini olusturur.
  \item WebView bu URL'i acar.
  \item Chart sayfasi kendi icinden ayni origin altindaki \texttt{/api/v1/candles}, \texttt{/api/v1/indicators}, \texttt{/api/v1/backtests/*} gibi rotalara istek atar.
  \item Bu rotalar da arka planda \texttt{evalon-backtest-api} Cloud Run servisine proxy olur.
\end{enumerate}

Bu sayede web veya mobile tarafinin ayrica candle API, indicator API ya da backtest API cagirisi yapmasi gerekmez.

\section{Desteklenen Query Parametreleri}

\begin{longtable}{|>{\ttfamily}p{3cm}|p{3cm}|p{7cm}|}
\hline
\textnormal{Parametre} & \textnormal{Tip} & \textnormal{Aciklama} \\
\hline
\textnormal{symbol} & string & Zorunlu hisse kodu. Ornek: \texttt{ADEL}, \texttt{THYAO}, \texttt{SISE}. \\
\hline
\textnormal{tf} & string & Zaman dilimi. Ornek: \texttt{1m}, \texttt{5m}, \texttt{15m}, \texttt{1h}, \texttt{4h}, \texttt{1d}, \texttt{1w}. \\
\hline
\end{longtable}

Ornek URL'ler:

\begin{lstlisting}[style=evalon]
/chart?symbol=THYAO&tf=1h
/chart?symbol=SISE&tf=4h
/chart?symbol=ADEL&tf=1d
\end{lstlisting}

\section{Web Tarafinda Nasil Kullanilir}

\subsection{Temel Mantik}

Web uygulamasi kullanicidan bir hisse ve timeframe secer. Sonra su sekilde URL olusturur:

\begin{lstlisting}[style=evalon]
const GRAPH_HOST =
  "https://evalon-graph-web-474112640179.europe-west1.run.app";

function buildChartUrl(symbol, tf = "1d") {
  const params = new URLSearchParams({
    symbol: symbol.toUpperCase(),
    tf,
  });
  return `${GRAPH_HOST}/chart?${params.toString()}`;
}
\end{lstlisting}

\subsection{Iframe veya Embedded WebView}

Web tarafinda tam ekran bir alan icine \texttt{iframe} ile koyabilirsiniz:

\begin{lstlisting}[style=evalon]
const url = buildChartUrl("ADEL", "1d");
iframe.src = url;
\end{lstlisting}

React / Next mantigi:

\begin{lstlisting}[style=evalon]
const src = buildChartUrl(selectedTicker, selectedTf);

return (
  <iframe
    src={src}
    style={{ width: "100%", height: "100vh", border: "0" }}
    allow="clipboard-read; clipboard-write"
  />
);
\end{lstlisting}

Not:

\begin{itemize}
  \item Web tarafinda tam ekran modal veya yeni route icinde iframe gostermek genelde en pratik yontemdir.
  \item Hisse degistiginde yeni URL uretilip iframe \texttt{src} degeri guncellenir.
\end{itemize}

\section{Mobile Tarafinda Nasil Kullanilir}

\subsection{Android WebView Ornegi}

\begin{lstlisting}[style=evalon]
val graphHost = "https://evalon-graph-web-474112640179.europe-west1.run.app"

fun buildChartUrl(symbol: String, tf: String = "1d"): String {
    return "$graphHost/chart?symbol=${symbol.uppercase()}&tf=$tf"
}

webView.settings.javaScriptEnabled = true
webView.settings.domStorageEnabled = true
webView.loadUrl(buildChartUrl("ADEL", "1d"))
\end{lstlisting}

\subsection{iOS WKWebView Ornegi}

\begin{lstlisting}[style=evalon]
let graphHost = "https://evalon-graph-web-474112640179.europe-west1.run.app"

func buildChartUrl(symbol: String, tf: String = "1d") -> URL {
    var comps = URLComponents(string: "\(graphHost)/chart")!
    comps.queryItems = [
        URLQueryItem(name: "symbol", value: symbol.uppercased()),
        URLQueryItem(name: "tf", value: tf)
    ]
    return comps.url!
}
\end{lstlisting}

Ardindan:

\begin{lstlisting}[style=evalon]
webView.load(URLRequest(url: buildChartUrl(symbol: "ADEL", tf: "1d")))
\end{lstlisting}

\subsection{Mobil Akis}

\begin{enumerate}
  \item Kullanici uygulamada hisseyi secer.
  \item Uygulama \texttt{buildChartUrl(selectedTicker, selectedTf)} ile URL olusturur.
  \item WebView bu URL ile acilir.
  \item Sayfa kendi icinde verileri yukler.
\end{enumerate}

\section{Bu Yapiyla Ne Kazanildi}

\begin{itemize}
  \item Chart sayfasi artik Vercel'e bagimli degil.
  \item Web tarafi yalnizca URL aciyor; API detayini chart hostunun kendisi yonetiyor.
  \item Mobile tarafi da ayni URL formatiyla ayni chart deneyimini kullanabiliyor.
  \item Ayni host uzerinden chart, backtest ve AI sayfalari ileride ayni mantikla WebView olarak kullanilabilir.
\end{itemize}

\section{Deploy Ozet Bilgisi}

Bu servis su ayarlarla deploy edildi:

\begin{itemize}
  \item Project: \texttt{evalon-490523}
  \item Region: \texttt{europe-west1}
  \item Service: \texttt{evalon-graph-web}
  \item Runtime account: \texttt{evalon-graph-web-run@evalon-490523.iam.gserviceaccount.com}
  \item Upstream backtest API: \texttt{https://evalon-backtest-api-474112640179.europe-west1.run.app/v1}
\end{itemize}

Repo icinde eklenen onemli dosyalar:

\begin{itemize}
  \item \texttt{backend/graph/Dockerfile}
  \item \texttt{backend/graph/scripts/deploy\_cloud\_run.sh}
  \item \texttt{backend/graph/scripts/smoke\_test\_cloud\_run.sh}
  \item \texttt{backend/graph/DEPLOY\_CLOUD\_RUN.md}
\end{itemize}

\section{Smoke Test Sonucu}

Canli host uzerinde su kontroller gecti:

\begin{itemize}
  \item \texttt{/health} 200
  \item \texttt{/chart.html?symbol=ADEL\&tf=1d} 200
  \item \texttt{/chart?symbol=ADEL\&tf=1d} 200
  \item \texttt{/api/v1/candles?symbol=ADEL\&tf=1d\&limit=2} 200
  \item \texttt{/api/v1/indicators?symbol=ADEL\&tf=1d\&strategy=rsi} 200
\end{itemize}

\section{Sonuc}

Artik web ve mobile uygulama tarafi su mantikla ilerleyebilir:

\begin{enumerate}
  \item Kullanici hisse secer.
  \item Uygulama Cloud Run chart URL'ini uretir.
  \item WebView bu URL'i acar.
  \item Chart, indikator ve diger veri isteklerini ayni host uzerinden kendi yonetir.
\end{enumerate}

Onerilen production kullanim formati:

\begin{lstlisting}[style=evalon]
https://evalon-graph-web-474112640179.europe-west1.run.app/chart?symbol=THYAO&tf=1d
\end{lstlisting}

\end{document}
